% Worked Examples: Concept 1.1.3 - State Space Models
\documentclass[11pt,a4paper]{article}
\usepackage{worked-examples-template}

\title{\textbf{Worked Examples}\\[0.5em]
\large Concept 1.1.3: State Space Models\\[0.3em]
\normalsize \coursesubtitle{} -- \coursetitle}
\author{Assoc.\ Prof.\ William Robertson \and Dr Sean McGowan}
\date{\today}

\begin{document}

\maketitle
\tableofcontents
\newpage

\section{Concept 1.1.3: State Space}

\subsection{Example 1: Converting RLC Circuit to State-Space Form}

\paragraph{Problem:} Consider an RLC circuit with:
\begin{itemize}
\item Resistor: $R = 10$ $\Omega$
\item Inductor: $L = 0.5$ H
\item Capacitor: $C = 0.01$ F
\item Input voltage: $v_{in}(t)$
\item Output: capacitor voltage $v_C(t)$
\end{itemize}

\begin{figure}[h]
\centering
\begin{circuitikz}[scale=1.0]
  % Input voltage source
  \draw (0,0) to[V=$v_{in}$] (0,3);
  
  % Resistor
  \draw (0,3) to[R=$R$, i=$i$] (2,3);
  
  % Inductor
  \draw (2,3) to[L=$L$] (4,3);
  
  % Capacitor and output
  \draw (4,3) to[C=$C$, v=$v_C$] (4,0);
  
  % Ground
  \draw (0,0) -- (4,0);
  \draw (4,0) node[ground] {};
  
  % Output label
  \draw[->] (4,3) -- (5,3) node[right] {$v_{out} = v_C$};
\end{circuitikz}
\caption{Series RLC circuit with voltage input and capacitor voltage output.}
\end{figure}

(a) Derive the differential equations governing the circuit.

(b) Choose appropriate state variables and write the system in state-space form $\dot{x} = Ax + Bu$, $y = Cx + Du$.

(c) Substitute the numerical values and write the state-space matrices explicitly.

(d) Determine the eigenvalues and characterize the system behavior.

\paragraph{Solution:}

\textbf{Step 1: Circuit differential equations}

Using Kirchhoff's voltage law around the loop:
\begin{align}
v_{in}(t) = Ri(t) + L\frac{di(t)}{dt} + v_C(t)
\end{align}

The current-voltage relationship for the capacitor:
\begin{align}
i(t) = C\frac{dv_C(t)}{dt}
\end{align}

\textbf{Step 2: State variable selection}

Choose state variables:
\begin{align}
x_1 &= v_C \quad \text{(capacitor voltage)} \\
x_2 &= i \quad \text{(inductor current)}
\end{align}

From the capacitor equation:
\begin{align}
\frac{dx_1}{dt} = \frac{dv_C}{dt} = \frac{i}{C} = \frac{x_2}{C}
\end{align}

From the KVL equation, solve for $\frac{di}{dt}$:
\begin{align}
L\frac{di}{dt} = v_{in} - Ri - v_C
\end{align}
\begin{align}
\frac{dx_2}{dt} = \frac{di}{dt} = \frac{v_{in}}{L} - \frac{R}{L}i - \frac{1}{L}v_C = -\frac{1}{L}x_1 - \frac{R}{L}x_2 + \frac{1}{L}v_{in}
\end{align}

\textbf{Step 3: State-space form}

The state equations are:
\begin{align}
\frac{d}{dt}\begin{bmatrix} x_1 \\ x_2 \end{bmatrix} = \begin{bmatrix} 0 & \frac{1}{C} \\ -\frac{1}{L} & -\frac{R}{L} \end{bmatrix}\begin{bmatrix} x_1 \\ x_2 \end{bmatrix} + \begin{bmatrix} 0 \\ \frac{1}{L} \end{bmatrix}v_{in}
\end{align}

The output equation (measuring capacitor voltage):
\begin{align}
y = v_C = \begin{bmatrix} 1 & 0 \end{bmatrix}\begin{bmatrix} x_1 \\ x_2 \end{bmatrix}
\end{align}

\textbf{Step 4: Numerical state-space matrices}

Substituting $R = 10$ $\Omega$, $L = 0.5$ H, $C = 0.01$ F:
\begin{align}
A &= \begin{bmatrix} 0 & \frac{1}{0.01} \\ -\frac{1}{0.5} & -\frac{10}{0.5} \end{bmatrix} = \begin{bmatrix} 0 & 100 \\ -2 & -20 \end{bmatrix} \\
B &= \begin{bmatrix} 0 \\ \frac{1}{0.5} \end{bmatrix} = \begin{bmatrix} 0 \\ 2 \end{bmatrix} \\
C &= \begin{bmatrix} 1 & 0 \end{bmatrix} \\
D &= 0
\end{align}

\textbf{Step 5: Eigenvalue analysis}

The characteristic equation is:
\begin{align}
\det(sI - A) = \det\begin{bmatrix} s & -100 \\ 2 & s+20 \end{bmatrix} = s(s+20) + 200 = s^2 + 20s + 200 = 0
\end{align}

Using the quadratic formula:
\begin{align}
s = \frac{-20 \pm \sqrt{400 - 800}}{2} = \frac{-20 \pm \sqrt{-400}}{2} = \frac{-20 \pm 20i}{2} = -10 \pm 10i
\end{align}

The eigenvalues are $\lambda_1 = -10 + 10i$ and $\lambda_2 = -10 - 10i$.

\textbf{System characterization:}
\begin{itemize}
\item \textbf{Damping ratio:} From the standard form $s^2 + 2\zeta\omega_n s + \omega_n^2$:
\begin{align}
\omega_n = \sqrt{200} = 14.14 \text{ rad/s}, \quad 2\zeta\omega_n = 20 \quad \Rightarrow \quad \zeta = \frac{20}{2 \times 14.14} = 0.707
\end{align}
\item \textbf{Behavior:} The system is \textbf{underdamped} with critical damping ratio $\zeta = 0.707$ (exactly at the boundary between underdamped and critically damped).
\item \textbf{Natural frequency:} $\omega_n = 14.14$ rad/s
\item \textbf{Stability:} Both eigenvalues have negative real parts, so the system is \textbf{stable}.
\item \textbf{Time constant:} $\tau = 1/\Re(\lambda) = 1/10 = 0.1$ s
\end{itemize}

\textbf{Key Insights:}
\begin{itemize}
\item State-space representation provides a systematic way to model dynamic systems
\item State variables represent the energy storage elements (capacitor voltage, inductor current)
\item The state-space form is particularly useful for multi-input, multi-output (MIMO) systems and modern control design
\item Eigenvalues of the $A$ matrix determine the stability and dynamic characteristics
\item This RLC circuit example is analogous to a mass-spring-damper system: $C \leftrightarrow m$, $L \leftrightarrow k$, $R \leftrightarrow c$
\end{itemize}
\end{document}
